\section{Kết luận và hướng phát triển}

\subsection{Kết luận}

Dự án \projectname{} đã hình thành một pipeline tương đối hoàn chỉnh từ camera/video đến kế hoạch đèn và bộ điều khiển nhúng. Điểm quan trọng nhất là hệ thống không dừng ở việc nhận diện xe. Phần AI chỉ tạo tín hiệu mức cao; MCU mới là thành phần thực thi điều khiển đèn theo thời gian thực.

Về mặt Vi xử lý, dự án đã thể hiện các nội dung cốt lõi: giao tiếp UART, parse và validate bản tin, queue giữa các task, task notification, mutex bảo vệ tài nguyên dùng chung, software timer, watchdog, FSM cục bộ và GPIO output. Các nội dung này phù hợp hơn với trọng tâm môn học so với việc chỉ trình bày độ chính xác của detector.

Về mặt thị giác máy tính, hệ thống đã có detector YOLO cục bộ, ROI cấu hình bằng JSON, đếm theo hướng, mật độ PCE, EMA và scheduler rule-based. Thiết kế này đủ giải thích, dễ kiểm thử và có thể mở rộng sang Raspberry Pi.

\subsection{Hướng phát triển}

Các bước tiếp theo nên được thực hiện theo thứ tự:

\begin{enumerate}
  \item Kiểm thử mạch đèn thật với ESP32, xác nhận GPIO, LED, nguồn và điện trở.
  \item Chạy hardware-in-the-loop đầy đủ: YOLO count \(\rightarrow\) signal plan \(\rightarrow\) UART \(\rightarrow\) ESP32 \(\rightarrow\) đèn thật.
  \item Chuẩn hóa thí nghiệm benchmark: ghi rõ phần cứng, model, video, config, commit, FPS, độ trễ UART, số lần timeout và số lần queue overflow.
  \item Tích hợp bảo mật vào runtime thật: HMAC cho PLAN, kiểm tra MAC phía MCU, chống replay và xác thực ACK.
  \item Tối ưu triển khai Raspberry Pi: camera thật, ROI camera thật, frame skipping, ONNX/NCNN nếu PyTorch không đủ FPS.
  \item Bổ sung mô phỏng hoặc dữ liệu đối chứng để đo thời gian chờ, hàng chờ và hiệu quả so với fixed-time controller.
\end{enumerate}

Khi các bước này hoàn tất, báo cáo có thể nâng cấp từ mức mô hình thử nghiệm sang mức đánh giá hệ thống nhúng AI có dữ liệu định lượng đầy đủ.

